% Subsection: What It Means for a Family to Be Closed  (notes stub; content authored later)
\subsection{What It Means for a Family to Be Closed}

\begin{remark*}[Exposition]
An operation on subsets always produces another subset of the ambient set.
Closure asks for more: it asks whether the result is still one of the selected
subsets.  Thus closure is an internal-stability condition on a family of
subsets.
\end{remark*}

\begin{definition}[Closed Under an Operation]\label{def:closed-under-operation}
Let \(X\) be a set, let \(\mathcal{F}\subseteq\mathcal{P}(X)\), and let
\(\Phi\) be an operation that takes one or more subsets of \(X\) as inputs and
returns a subset of \(X\).  The family \(\mathcal{F}\) is \emph{closed under}
\(\Phi\) if every admissible choice of inputs from \(\mathcal{F}\) has output
again in \(\mathcal{F}\).
\end{definition}

\begin{remark*}[Standard quantified statement]
For a binary operation \(\Phi\), closure under \(\Phi\) means
\[
  \forall A,B\in\mathcal{F},\qquad
  \Phi(A,B)\in\mathcal{F}.
\]
For a unary operation \(\Psi\), closure under \(\Psi\) means
\[
  \forall A\in\mathcal{F},\qquad
  \Psi(A)\in\mathcal{F}.
\]
\end{remark*}

\begin{remark*}[Predicate reading]
\[
\begin{aligned}
&C_{\Theta}=\mathsf{ClosureContext}(X,D_{\Theta},\mathcal{F},\Theta),\\
&\operatorname{ClosedUnder}(\Theta,\mathcal{F},C_{\Theta}).
\end{aligned}
\]
\end{remark*}

\begin{remark*}[Interpretation]
Closure is not a statement about whether an operation can be performed.  It is
a statement about whether performing the operation stays inside the chosen
family.  The ambient power set \(\mathcal{P}(X)\) is closed under all ordinary
set operations, but a smaller family may fail closure for some of them.
\end{remark*}

\begin{remark*}[Examples]
If \(\mathcal{F}\) is the family of finite subsets of \(X\), then
\(\mathcal{F}\) is closed under finite unions.  If \(X\) is infinite, the same
family is not closed under complements, because the complement of a finite
subset may be infinite.
\end{remark*}

\begin{dependencies}
\begin{itemize}
  \item \hyperref[def:family-of-subsets]{Family of Subsets}
  \item \hyperref[def:operation-on-members-of-family]{Operation on Members of a Family}
  \item \hyperref[def:power-set]{Power Set}
\end{itemize}
\end{dependencies}
